\section{Discussion}
\label{sec:discussion}


% \begin{figure}[t!]
%     \begin{tcolorbox}[breakable,leftrule=1.5mm,top=0.8mm,bottom=0.5mm,colback=blue!5,colframe=blue!75!black]
%     \textbf{Takeaway \#1: Think through your metrics!}\\
%     If one cares about scaling-predictable evaluations, then one needs to think through how their evaluations transform raw model outputs into useful signals to know what to expect.\\[0.6em]
%     \textbf{Takeaway \#2: Continuous metrics are insufficient to guarantee predictable changes.}\\
%     As shown by $p_{\theta}^{\text{Choices}}(\text{Correct Choice})$ \& \texttt{Brier Score}, even ``continuous'' metrics can be unpredictable,
%     e.g., if the metric weighs correct behavior against specific incorrect behaviors.\\[0.6em]
%     %
%     \textbf{Takeaway \#3: Recommended scaling-predictable metrics for pretraining practitioners.}\\
%     Pretraining practitioners seeking scaling-predictable signals for capabilities are advised to focus on \(p_{\theta}^\textit{Vocab}(\text{Correct Choice})\) on relevant benchmarks. Scores under this metric provide smoother scaling trends and are arguably more interpretable than the pretraining loss.\\[0.6em]
%     \textbf{Takeaway \#4: Evaluations should be reshaped based on intended desiderata.}\\
%     Too often, we take evaluations as frozen static objects, but evaluations should be adapted to pertinent goals. For instance, if the goal is to predict capabilities with scale, evaluations should be designed or adapted to be scaling-predictable.
%     \end{tcolorbox}
%     \label{takeaways}
% \end{figure}
This work identifies why multiple-choice assessments of frontier AI models are unpredictable, as well as the underlying mechanism: probability mass on incorrect choices.
%
Our results have implications for the design of future evaluations of frontier AI models that are reliably predictable with scaling. We hope that our work will be extended to further the science of scaling-predictable evaluation of AI systems, especially for complex and important model capabilities.
We note several future directions for extension of our work, and we hope that the community also adopts our framing to further improve scaling-predictable evaluations. 
%
For \textbf{Related Work}, please see Appendix ~\ref{app:sec:related_work}.
For \textbf{Future Directions}, please see Appendix~\ref{app:sec:future_directions}.



% \section*{Acknowledgements}


% We thank a number of colleagues for their feedback in improving this work: Tamay Besiroglu, Cl\'ementine Fourrier, Geoffrey Irving, David Owen, Lintang Sutawika and Joshua Kazdan. R.S. acknowledges support from Stanford Data Science and from OpenAI's Superalignment Fast Grant Research Fellowship. S.K. is partially supported by NSF III 2046795, IIS 1909577, CCF 1934986, NIH 1R01MH116226-01A, NIFA award 2020-67021-32799, the Alfred P. Sloan Foundation, and Google Inc.